% Part: computability
% Chapter: computability-theory
% Section: normal-form

\documentclass[../../../include/open-logic-section]{subfiles}

\begin{document}

\olfileid[fr]{cmp}{thy}{nfm}
\olsection{Le théorème de forme normale}

Supposons que nous puissions décrire les définitions des
fonctions calculables et vérifier si un enregistrement
proposé constitue bien le relevé complet du calcul d'une
de ces fonctions sur une entrée donnée. Nous pouvons
alors déterminer, quel que soit le modèle de calcul
usuel considéré, la valeur d'une fonction calculable~$f$
sur une entrée~$x$ en procédant ainsi :
\begin{enumerate}
  \item Parcourir toutes les descriptions possibles
  d'enregistrements de calcul.
  \item Vérifier si l'un de ces enregistrements décrit
  le calcul de~$f(x)$.
  \item Si un enregistrement correct est trouvé,
  en extraire la valeur de~$f(x)$.
\end{enumerate}
Le théorème de forme normale de Kleene donne un
fondement rigoureux à cette procédure.

\begin{thm}[Théorème de forme normale de Kleene]
\ollabel{thm:normal-form}
Il existe une relation récursive primitive~$T(e,x,s)$
et une fonction récursive primitive~$U(s)$ ayant la
propriété suivante : pour toute fonction partielle
calculable \emph{unaire}~$f$, il existe un entier~$e$ tel que
\[
f(x) \simeq U(\umin{s}{T(e, x, s)})
\]
pour tout~$x$. Les fonctions de plusieurs arguments
se traitent en codant leurs tuples d'arguments ou
en adaptant~$T$ à chaque arité.\footnote{Précisions éditoriales : nous explicitons l'arité unaire de la formule à un argument de la source. Les affirmations sur les modèles de calcul sont comprises dans le cadre des modèles effectifs usuels et d'un codage approprié.}
\end{thm}

\begin{proof}[Esquisse de démonstration]
Dans un modèle usuel de calcul, on peut définir
rigoureusement une description de la fonction
calculable~$f$ et coder cette description par un
entier naturel~$e$. On peut également définir une
« suite de calcul » qui enregistre le calcul de la
fonction d'indice~$e$ sur l'entrée~$x$, puis coder
cette suite par un entier~$s$. On choisit les codages
de telle sorte que
\begin{enumerate}
  \item la relation $T(e,x,s)$, vraie si et seulement
  si $s$ code la suite complète d'un calcul terminé de la fonction
  d'indice~$e$ sur l'entrée~$x$, et
  \item la fonction $U(s)$, qui associe à une suite
  de calcul codée par~$s$ le résultat final
  de ce calcul,
\end{enumerate}
soient toutes deux calculables. Avec un codage
approprié, la relation~$T$ et la fonction~$U$
sont même récursives primitives.
\end{proof}

\begin{explain}
Pour démontrer rigoureusement le théorème de forme
normale, il faudrait fixer un modèle de calcul,
coder en détail les descriptions de fonctions
calculables et les suites de calcul, puis vérifier
que la relation~$T$ et la fonction~$U$ sont
récursives primitives. Dans la plupart des
applications, il suffit que $T$ et~$U$ soient
calculables et que $U$ soit totale. On peut, par convention, lui donner la valeur $0$ sur les entiers qui ne codent pas un calcul terminé. Cette extension ne change aucune valeur obtenue à partir d'un témoin de $T$.

On peut retenir la preuve sous la forme suivante :
$\umin{s}{T(e, x, s)}$ cherche une suite de calcul
de la fonction d'indice~$e$ sur l'entrée~$x$,
et $U$ renvoie le résultat de ce calcul lorsqu'une
telle suite est trouvée.
\end{explain}

Lorsque le modèle de calcul est celui des fonctions
récursives partielles, employé à l'origine par
Kleene, le théorème montre qu'une seule application
de la recherche non bornée, c'est-à-dire d'un
opérateur $\umin{y}{f(\vec x,y)}$, suffit à définir
toute fonction récursive partielle. C'est en ce sens
qu'une telle fonction possède une forme normale.

$T$ et $U$ permettent de définir l'énumération
$\cfind{0}$, $\cfind{1}$, $\cfind{2}$, \dots.
Désormais, nous supposerons fixé un choix approprié
de $T$ et de~$U$ pour le codage usuel des programmes, et nous prendrons l'équation
\[
\cfind{e}(x) \simeq U(\umin{s}{T(e,x,s)})
\]
comme \emph{définition} de $\cfind{e}$.

Voici encore un fait utile :

\begin{thm}
Dans le codage usuel des programmes, toute fonction partielle calculable possède une
infinité d'indices.
\end{thm}

Cela se comprend intuitivement. Partant d'une
description d'une fonction calculable, on peut
en produire une autre qui calcule la même
fonction, c'est-à-dire les mêmes couples
entrée-résultat, après avoir effectué une
opération sans effet sur le résultat, par exemple
tester si $0=0$ ou compter jusqu'à~$5$.
La nouvelle description possède un indice
différent de l'originale. Les deux indices
désignent la même fonction calculée de deux
manières légèrement différentes. En ajoutant un nombre arbitraire de telles opérations inutiles, on obtient une infinité de descriptions distinctes, donc d'indices distincts dans ce codage.

\end{document}
